\documentclass[11pt]{article}
\usepackage[margin=0.85in]{geometry}
\usepackage{amsmath,amssymb}
\usepackage[hidelinks]{hyperref}
\usepackage{enumitem}

\setlength{\parindent}{0pt}
\setlength{\parskip}{0.55\baselineskip}
\setlength{\emergencystretch}{3em}
\setlist{nosep,leftmargin=1.5em}
\urlstyle{same}

\title{Technical Review of\\
\emph{Public Emergency Codes: Evidence Estimate and Full-Implementation Potential}}
\author{}
\date{August 8, 2026}

\begin{document}
\maketitle

\section*{Overall assessment}

The manuscript has a useful high-level distinction between a current-evidence
estimand and a mature-implementation counterfactual, and the reported
deterministic arithmetic is internally correct.  Running
\texttt{pec\_evidence\_simulation.py} and
\texttt{pec\_full\_potential\_simulation.py} with the stated seeds reproduces
the numerical tables.  Equations (2) and (3) are dimensionally consistent if
$k$ is interpreted exactly as a relative survival change per minute.  The
reported per-million rates and the sum of the pathway values in Table 7 also
check arithmetically.

The paper nevertheless requires major revision before its headline numbers can
be described as evidence-supported estimates.  The principal problem is not
arithmetic; it is causal identification.  Every positive component of the
$66.88$-life modal calculation depends on at least one judgmental causal or
PEC-attribution parameter that is not estimated by the cited studies.  The
$2{,}287$ full-implementation value relies on a second, much larger set of
unsourced absolute mortality-risk changes.  These calculations can be retained
as transparent planning scenarios, but the present labels overstate what the
evidence identifies.

No inverse-trigonometric expressions, coordinate transformations, or branch
choices occur in the manuscript, so no branch or quadrant defect was found.
The current-model sign convention is clear.  The full-potential notation and
the treatment of overlapping pathways remain implementation-ambiguous, as
detailed below.

\section*{Highest-priority technical findings}

\subsection*{1. The ``current-evidence'' point estimate is not identified by current PEC evidence}

\textbf{Classification: unsupported claim.}

\textbf{Location.} Abstract; Sections 1, 5.1, 7, and 8; Equation (1).

\textbf{Problem.} The modal 988 contribution is
\[
L_{988}=\frac{4{,}372}{2.5}(0.04)(0.50)=34.98,
\]
which is approximately $52\%$ of the headline $66.88$ lives/year.  The source
study estimates an association for the entire national 988 rollout, whereas
the PEC share $g$ and causal-attribution factor $c$ are elicited planning
assumptions.  The manuscript explicitly acknowledges that there is no direct
PEC trial, but then describes the resulting positive value as supported by
current direct evidence.  Neither the $4\%$ mode for $g$ nor the $50\%$ mode
for $c$ is estimated from the cited awareness/use survey.

The other two positive modules also require the assumed slope $k$.  Thus the
full $66.88$-life modal result is a scenario conditional on causal assumptions,
not an evidence-identified PEC effect.

\textbf{Why it matters.} The central interpretation of the paper depends on a
claimed separation between evidence-supported lives and speculative potential.
The present construction does not maintain that separation.

\textbf{Suggested fix.} Report at least three layers separately:
\begin{enumerate}
\item observed external associations and measured process delays;
\item a PEC-attribution scenario conditional on $g$, $c$, and $k$; and
\item the mature full-implementation scenario.
\end{enumerate}
Rename $66.88$ the ``modal assumption-conditional current-deployment
scenario,'' or set the evidence-only PEC mortality contribution to zero until
PEC-specific attribution is measured.  If a positive current scenario is
retained, provide an explicit elicitation method and sensitivity analysis for
$g$ and $c$ rather than calling their ranges evidence-calibrated.

\subsection*{2. The cited evidence does not establish the modeled OHCA survival slope $k$}

\textbf{Classification: definite citation mismatch and unsupported causal
conversion.}

\textbf{Location.} Evidence-audit row ``Causal conversion to survival,''
Table 4, Equations (2)--(3), and bibliography items
\texttt{timecpr2026} and \texttt{responsetime2011}.

\textbf{Problem.} The paper applies
\[
k=3/5/8\%\quad\text{relative survival gain per minute}
\]
to the national OHCA population.  The DOI given for
\texttt{timecpr2026}, DOI
\url{https://doi.org/10.1161/CIRCOUTCOMES.125.013166}, is a pediatric
OHCA study: \emph{Time to Bystander Cardiopulmonary Resuscitation and Survival
Outcomes in Pediatric Out-of-Hospital Cardiac Arrests in the United States}.
The bibliography omits ``Pediatric'' and the rest of the actual title.  It
cannot by itself justify applying one common slope to the overwhelmingly adult
national OHCA population.

The study cited as \texttt{responsetime2011} examined all high-priority urban
EMS events rather than an OHCA-specific per-minute effect.  It reported a
$0.7$ percentage-point mortality difference across an eight-minute threshold,
with a $95\%$ confidence interval from $-0.5$ to $2.0$ percentage points and an
adjusted odds-ratio confidence interval including one.  It therefore does not
estimate the paper's positive $3$--$8\%$ relative gain per minute.

\textbf{Why it matters.} $k$ multiplies both OHCA modules and is shared across
their draws.  A population or estimand mismatch propagates directly into the
language and access life estimates and their apparent precision.

\textbf{Suggested fix.} Use an adult, OHCA-specific time-to-survival model and
show the exact conversion from its reported estimand to a relative change per
minute.  Prefer separate curves for time to CPR and time to patient contact;
they are not necessarily exchangeable.  Stratify pediatric and adult cases if
the pediatric source is retained.  Until an appropriate slope is justified,
move these mortality components to a clearly labeled scenario or zero-centered
sensitivity analysis.  Correct the bibliography title in all cases.

The source identities can be checked at
\url{https://www.ahajournals.org/doi/10.1161/CIRCOUTCOMES.125.013166}
and \url{https://pubmed.ncbi.nlm.nih.gov/22026820/}.

\subsection*{3. The full-implementation mortality effects are planning priors, not evidence-calibrated effects}

\textbf{Classification: unsupported claim and reproducibility gap in the
paper text.}

\textbf{Location.} Section 9, Equation (7), Tables 7--8, and
\texttt{pec\_full\_potential\_simulation.py}.

\textbf{Problem.} The script contains low/mode/high distributions for
$r_j$, beneficial absolute risk change, adverse absolute risk change,
end-to-end success, and overlap for nine pathways.  Most of these distributions
do not appear in the manuscript, and no source or elicitation procedure is
given for the absolute mortality-risk changes.  Examples include the positive
means assigned to 211 referral, 911 diversion, medical data/photo/video, and
passive monitoring even though the evidence audit states that these pathways
lack demonstrated mortality effects.  Mature implementation can increase
delivery success, but it does not establish a clinical effect that is absent
from the evidence.

\textbf{Why it matters.} Table 7 exposes only the modal relevance, high
success, net modal risk change, and high attribution used for one plug-in
calculation.  A reader cannot reconstruct the $90\%$ interval from the paper,
assess the plausibility of the tails, or determine which assumptions create the
$3{,}054$ median without reading source code.  More importantly, the numerical
range may be mistaken for empirically calibrated uncertainty.

\textbf{Suggested fix.} Add a complete full-potential input table containing
$I_j$, every low/mode/high tuple, parameter units, evidence tier, source, and
elicitation rationale.  Mark unsupported mortality changes as expert-judgment
priors.  Report a zero-mortality-effect version for Tier C pathways alongside
the planning scenario.  Archive the exact script and a machine-readable input
file with the paper.

\subsection*{4. Scalar attribution factors do not resolve overlap or shared uncertainty}

\textbf{Classification: likely upward point bias and likely understated
aggregate uncertainty.}

\textbf{Location.} Evidence-audit row ``Feature overlap and correlation,''
Equation (7), Table 7, and the full-potential script.

\textbf{Problem.} Several full-potential rows use the same
$36.367261$-million treated-and-transported 911 denominator.  A single incident
can simultaneously involve earlier 911 access, language support, location
assistance, medical data, video, and contact alerts.  Multiplying each row by a
separate $w_{F,j}$ does not make the rows mutually exclusive and does not show
which pathway receives the shared outcome.

The script also samples the major pathway inputs independently.  National
adoption, device reach, platform reliability, PSAP compatibility, and causal
credibility would create strong positive correlations across pathways.
Independent sampling can make the aggregate interval artificially narrow even
when every marginal range is wide.

The current model similarly calls the language and access mechanisms
``non-overlapping,'' although they can occur in the same OHCA.  Additivity may
be defensible if the recovered times are sequential and the survival response
is locally linear, but that assumption is not stated or tested.

\textbf{Why it matters.} The headline full-potential result is a sum across
pathways.  Double counting affects the point estimate, while omitted shared
uncertainty affects the interval and threshold probabilities.

\textbf{Suggested fix.} Use an incident-level model with mutually exclusive
clinical states and explicit combinations of enabled functions.  Alternatively,
define a formal allocation rule for overlapping risk reductions and show that
the total change in an incident cannot exceed its baseline mortality risk.
Introduce shared latent draws for national reach, user activation, technical
success, and receiver compatibility, then report a dependence sensitivity
analysis.

\subsection*{5. The reported point summaries and interval terminology are internally inconsistent}

\textbf{Classification: definite reporting inconsistency.}

\textbf{Location.} Abstract and Tables 6--8.

\textbf{Problem.} The manuscript calls $2{,}287$ the ``central''
full-implementation potential, but its own simulation gives
\[
\Pr(L_F>2{,}287)=86.4\%.
\]
Thus the plug-in value is near the lower tail, not the center of the simulated
distribution; the mean is $3{,}148$ and the median is $3{,}054$.  The analogous
current calculation headlines the modal plug-in value $66.9$, whereas the
predictive mean and median are approximately $75.7$ and $75.3$.

Table 8 also labels the full-potential mean and median ``predictive,'' although
the full-potential script adds parameter/scenario uncertainty only and does not
simulate annual outcome noise.  By contrast, the current-evidence script does
add an explicit annual-noise term.

\textbf{Why it matters.} Readers may interpret $2{,}287$ as a median or
expected value and the full interval as a predictive interval for future
observed outcomes.  Neither interpretation matches the implementation.

\textbf{Suggested fix.} Rename $2{,}287$ the ``joint modal-input plug-in
scenario'' and use either the median or mean as the probabilistic headline.
Use ``parameter-uncertainty interval'' or ``scenario interval'' for the full
model unless outcome noise is actually added.  Apply the same point-summary
convention to both estimands and state why that summary is decision-relevant.

\subsection*{6. The benefit probabilities are dominated by subjective zero-centered budgets}

\textbf{Classification: implementation-facing ambiguity and overstated
precision.}

\textbf{Location.} Section 5.4, Table 5, and Section 8.

\textbf{Problem.} The six independent normal budgets have a combined standard
deviation of
\[
\sqrt{35^2+12^2+10^2+8^2+18^2+25^2}=49.82
\quad\text{lives/year}.
\]
In the reproduced simulation, the direct positive modules have a $5$th--$95$th
percentile range of approximately $44.6$--$113.0$, whereas adding the
zero-centered budgets expands it to approximately $-12.7$--$165.2$ before
annual outcome noise.  Consequently, quantities such as
$\Pr(L_{\mathrm{net}}>0)=91.7\%$ depend strongly on six unvalidated standard
deviations and their assumed independence.

The added annual noise is described only as a normal approximation with
standard deviation $\sqrt{\max(L_{\mathrm{direct}},1)}$.  No outcome-generating
argument is given for that form, and negative contributions preclude a direct
Poisson interpretation.

\textbf{Why it matters.} The threshold probabilities can look like inferential
probabilities derived from data even though they are primarily conditional on
planning-budget choices.

\textbf{Suggested fix.} Label these outputs prior-predictive or
assumption-conditional.  Report results with the budgets removed and with all
$\sigma_m$ multiplied by, for example, $0.5$, $1$, and $2$.  Justify their
dependence structure.  Either derive the annual-noise model from explicit event
counts or omit it and call the result a parameter-uncertainty distribution.

\subsection*{7. The national OHCA expansion needs a defensible denominator model}

\textbf{Classification: likely issue requiring external verification.}

\textbf{Location.} Section 2 and Table 4.

\textbf{Problem.} The calculation $137{,}119/0.56\approx245{,}000$ treats the
rounded CARES population-coverage percentage as though it were an exact case
sampling fraction with nationally representative incidence.  CARES describes a
catchment population, not a probability sample of U.S. OHCA events.  Geographic
differences in age, incidence, EMS participation, and case capture can make a
simple inverse-coverage expansion biased.

\textbf{Why it matters.} Both positive OHCA modules scale linearly with this
anchor, and the same anchor is used in the full-potential model.

\textbf{Suggested fix.} Use the exact covered population and a documented
age-/sex-standardized national incidence estimate, or retain the expansion but
give it an uncertainty distribution and sensitivity analysis.  State explicitly
whether the numerator and coverage percentage use identical case definitions
and dates.

\section*{Notation and equation audit}

\begin{itemize}
\item $L_E$ is defined in Section 1 but is not subsequently used; the paper
switches to $L_{\mathrm{net}}$ and $L_{\mathrm{net,base}}$.  Define
$L_E\equiv L_{\mathrm{net}}$ for the current-evidence model or use one symbol
throughout.
\item In Equation (7), $\Delta_j$ and $h_j$ require explicit units and verbal
definitions.  Table 7 instead uses ``Net $\Delta p$,'' while the code uses
separate beneficial and adverse absolute risk changes.  Use one notation and
state that $\Delta_j-h_j$ is an absolute probability change per relevant,
successfully implemented incident.
\item The descriptors for $N_{\mathrm{OHCA}}$, $\pi_{\mathrm{LEP}}$,
$u_{\mathrm{LEP}}$, $\pi_A$, and $u_A$ are locally consistent.  However,
$\pi_{\mathrm{LEP}}$ is transferred from sampled medical 911 calls to OHCA and
$\pi_A$ is a planning sensitivity range.  Mark both as transported assumptions,
not OHCA-specific prevalence estimates.
\item Equations (2)--(3) have correct dimensions only under the relative-slope
interpretation of $k$.  If a source reports an absolute percentage-point change,
the equations and numerical inputs must be changed rather than inserting that
estimate directly as $k$.
\end{itemize}

\section*{Meaningful editorial and statistical-language issues}

\begin{itemize}
\item Table 2 says randomized video evidence ``found no mortality difference.''
Failure to reject a difference is not evidence of equivalence.  Replace this
with ``did not detect a statistically significant difference in 30-day
mortality'' and report the estimate and confidence interval if mortality is
used to justify a zero-centered prior.
\item Replace repeated phrases such as ``the current-evidence estimate is about
$67$'' with ``the modal plug-in scenario is about $67$'' unless the causal
parameters are empirically identified.
\item Table 1 places ``federal estimate'' in the ``Data year'' column for 911
calls.  Give the source publication year and access date, and distinguish a
current count from an older planning estimate.
\item The supplied-workspace documents cited as \texttt{pecinputs} and
\texttt{pecfunctions} are load-bearing sources for the full-potential
assumptions.  They should be included in the reproducibility archive or their
relevant content should be reproduced in the paper.
\end{itemize}

\section*{Proofing sweep}

The required mechanical scan was run on the compiled manuscript and its hits
were spot-checked.  The following high-confidence corrections remain:
\begin{itemize}
\item \textbf{Bibliography, \texttt{samhsa988}:} the identical ``Source page''
hyperlink is printed twice.  Delete one copy.
\item \textbf{Bibliography, \texttt{timecpr2026}:} restore the complete title,
including ``Survival Outcomes in Pediatric Out-of-Hospital Cardiac Arrests in
the United States,'' and add the authors, volume/issue or online-publication
details, and DOI in normal citation form.
\item \textbf{Table 2:} change ``found no mortality difference'' to the
non-equivalence wording recommended above.
\end{itemize}

No duplicated comma/period artifact, malformed ``into ... into'' derivation
frame, product-name capitalization error, or \texttt{arctan(x/y)} branch pattern
was confirmed.  The scan's semicolon-capitalization candidate in the Table 7
caption is grammatical because ``Function 9'' is a numbered proper label.

\section*{Highest-priority revision sequence}

\begin{enumerate}
\item Reclassify $66.88$ as an assumption-conditional scenario, or remove
positive PEC mortality attribution that is not empirically identified.
\item Replace or stratify the evidence for $k$ and show the exact causal
conversion from source estimates to the model parameter.
\item Publish every full-potential input and its evidentiary or elicitation
basis; add a Tier-C-zero scenario.
\item Rebuild overlap and dependence handling at the incident level or provide
a formal allocation and shared-factor sensitivity analysis.
\item Make the headline summary, interval terminology, and simulation code use
the same estimand definitions.
\item Correct the bibliography and statistical wording, then rerun both scripts
and update the abstract and quoted-result boxes from the revised model.
\end{enumerate}

\end{document}